\section{Toy Example for the Merge Case}

In this section, we formally define the $\mcase$. We first illustrate the main ideas through a simple toy example of the $\mcase$, and then generalize the result to the full setting.

In the $\mcase$, we assume that every cluster in the optimal solution at time step $t+1$, denoted by $\m{C}^{(t+1)^*}$, is either identical to a cluster in $\m{C}^{t^*}$ or is obtained by taking the union of a collection of clusters from $\m{C}^{t^*}$. We define it below formally.

\begin{definition}[$\mcase$]
We say that the sequence of optimal clusterings satisfies the $\mcase$ if for every cluster
$C^{t + 1} \in \m{C}^{(t+1)^*}$, one of the following holds:
\begin{enumerate}
    \item $C^{t + 1} \in \m{C}^{t^*}$, or
    \item there exists a collection of clusters
    \[
    C_1^{t}, C_2^{t}, \ldots, C_k^{t} \in \m{C}^{t^*}
    \]
    such that
    \[
    C = \bigcup_{i=1}^k C_i^{t} \cup \{v_{t+1}\}.
    \]
\end{enumerate}
\tienlong{"such that ... or $C = \bigcup_{i=1}^k C_i^{t} \cup \{v_{t+1}\}$"}
\end{definition}

Now, let us consider a toy example of the $\mcase$. Suppose that at time step $t_0$, the optimal solution $\optcls{0}$ consists of exactly two clusters, denoted by $S_1$ and $S_2$. Further, for every time step $t_k$, where $k \in \{1,\ldots,n\}$, suppose that the optimal solution $\optcls{k}$ consists of a single cluster
\[
    S = S_1 \cup S_2 \cup \{v_1,\ldots,v_k\},
\]
where $\{v_1,\ldots,v_k\}$ denotes the set of vertices that arrive at time steps $t_1,\ldots,t_k$, respectively.

We now describe the algorithm $\oncon$, which maintains a clustering $\solclsw$ at time step $t$. To define the algorithm, we first introduce the cost of assigning a vertex $v$ to an arbitrary subset $P \subseteq V$:
\[
    \cost(P,v)
    =
    \sum_{u \in P} w^{-}(u,v)
    +
    \sum_{u \in V \setminus P} w^{+}(u,v),
\]
where $w^{-}(u,v)$ denotes the number of input clusterings in which $u$ and $v$ are separated, and $w^{+}(u,v)$ denotes the number of input clusterings in which $u$ and $v$ are clustered together.

For every time step $t_k$, let $S_1^{t_k}$ and $S_2^{t_k}$ denote the current states of the clusters $S_1$ and $S_2$, respectively. We are now ready to describe the algorithm $\oncon$.

\paragraph{Algorithm $\oncon$.}
At time step $t_{k+1}$, upon the arrival of the vertex $v_{k+1}$, the algorithm performs the following update:
\begin{enumerate}[(i)]
    \item If $\cost(S_1^{t_k},v_{k+1}) \le \cost(S_2^{t_k},v_{k+1})$, then set
    \[
        \solcls{k+1}
        =
        \left\{
        S_1^{t_k} \cup \{v_{k+1}\},
        \,
        S_2^{t_k}
        \right\}.
    \]

    \item Otherwise, set
    \[
        \solcls{k+1}
        =
        \left\{
        S_1^{t_k},
        \,
        S_2^{t_k} \cup \{v_{k+1}\}
        \right\}.
    \]
\end{enumerate}
\tienlong{So we are not considering singleton cases.}
\kushagra{No}

We now prove the following lemma.

\begin{lemma}\label{lem:toy.example.merge.case}
If $\obj(\solcls{k}) \ge 2\,\obj(\optcls{0})$, then $\obj(\optcls{k}) \ge (1+\epsilon)\,\obj(\optcls{0})$ for some $\epsilon \in (0,1)$.
\end{lemma}

\begin{lemma}[Alternative charging proof for $s(N_1,N_2)$]\label{lem:charging.s.N1.N2}
Let
\[
    A=C_1,\qquad B=C_2,\qquad X=N_1,\qquad Y=N_2.
\]
Let
\[
    d(u,v)=\frac{w^-(u,v)}{m},
    \qquad
    s(u,v)=\frac{w^+(u,v)}{m}=1-d(u,v).
\]
where $w^-(u,v)$ is the number of clusterings where $u$ and $v$ are separated and $w^+(u,v)$ is the number of clusterings where $u$ and $v$ are together.
Assume that the vertices in $X\cup Y$ are assigned online and greedily to one
of the two current sides
\[
    A\cup X
    \qquad\text{and}\qquad
    B\cup Y.
\]
Assume also that
\[
    |X|+|Y|\le |A|+|B|.
\]
Define
\[
    I^{\OPT}
    =
    d(X,A)+d(X,B)+d(Y,A)+d(Y,B)+d(X,Y)+d(X)+d(Y).
\]
Suppose that for some constant $c>0$,
\[
    s(X,B)+d(X,A)+d(Y,B)+s(Y,A)
    \le
    c\,I^{\OPT}.
\]
Then
\[
    s(X,Y)\le K I^{\OPT},
\]
where $K$ is an absolute constant depending only on $c$. More precisely, for
any fixed $\alpha>2$,
\[
    s(X,Y)
    \le
    \left(
        \alpha+\frac{\alpha}{\alpha-1}+2c
    \right)I^{\OPT}.
\]
\end{lemma}

\begin{proof}
Let
\[
    a=|A|,
    \qquad
    b=|B|.
\]
For a vertex $v$ arriving online, let $X_{<v}$ and $Y_{<v}$ denote the vertices
of $X$ and $Y$, respectively, that arrived before $v$.

We split the arrivals into two cases. Fix a constant $\alpha>2$.

\paragraph{Low case.}
Suppose first that a vertex $v\in X$ is assigned to the $A$-side and satisfies
\[
    s(v,Y_{<v})
    \le
    \alpha\, d(v,B\cup Y_{<v}).
\]
Then
\[
    s(v,Y_{<v})
    \le
    \alpha d(v,B)+\alpha d(v,Y_{<v}).
\]
Similarly, if a vertex $v\in Y$ is assigned to the $B$-side and satisfies
\[
    s(v,X_{<v})
    \le
    \alpha\, d(v,A\cup X_{<v}),
\]
then
\[
    s(v,X_{<v})
    \le
    \alpha d(v,A)+\alpha d(v,X_{<v}).
\]

Summing over all low-case arrivals gives
\[
    \sum_{\text{low arrivals}} s(v,\text{opposite previous side})
    \le
    \alpha\bigl(d(X,B)+d(Y,A)+d(X,Y)\bigr).
\]
Since all three terms on the right are contained in $I^{\OPT}$, we obtain
\[
    \sum_{\text{low arrivals}} s(v,\text{opposite previous side})
    \le
    \alpha I^{\OPT}.
\]

\paragraph{High case into $X$.}
Now consider a vertex $v\in X$ assigned to the $A$-side. We call it high if
\[
    s(v,Y_{<v})
    >
    \alpha\, d(v,B\cup Y_{<v}).
\]
We first show that high arrivals into $X$ can occur only $O(|A|)$ times.

Let
\[
    x=|X_{<v}|,
    \qquad
    y=|Y_{<v}|.
\]
Define the gain of assigning $v$ to a set $S$ by
\[
    G(v,S)=s(v,S)-d(v,S).
\]
Since $s(u,v)+d(u,v)=1$ for every pair,
\[
    G(v,S)=|S|-2d(v,S).
\]
Because the algorithm assigns $v$ greedily to $A\cup X_{<v}$ instead of
$B\cup Y_{<v}$, we have
\[
    G(v,A\cup X_{<v})
    \ge
    G(v,B\cup Y_{<v}).
\]
Also,
\[
    G(v,A\cup X_{<v})\le a+x.
\]
Since $v$ is high,
\[
    s(v,Y_{<v})>\alpha\,d(v,B\cup Y_{<v}).
\]
Using the trivial bound $s(v,Y_{<v})\le y$, we get
\[
    d(v,B\cup Y_{<v})<\frac{y}{\alpha}.
\]
Therefore,
\[
    G(v,B\cup Y_{<v})
    =
    b+y-2d(v,B\cup Y_{<v})
    >
    b+\left(1-\frac{2}{\alpha}\right)y.
\]
Combining the last inequalities gives
\[
    a+x
    >
    b+\left(1-\frac{2}{\alpha}\right)y.
\]
Thus every high arrival into $X$ satisfies
\[
    x
    >
    b-a+\left(1-\frac{2}{\alpha}\right)y.
\]

Let $v^*$ be the first high arrival into $X$, and let
\[
    x^*=|X_{<v^*}|,
    \qquad
    y^*=|Y_{<v^*}|.
\]
Then
\[
    x^*
    >
    b-a+\left(1-\frac{2}{\alpha}\right)y^*.
\]
Let $H_X$ denote the number of high arrivals into $X$. Since all high arrivals
into $X$ occur from time $v^*$ onward,
\[
    H_X\le |X|-x^*.
\]
Using the size condition
\[
    |X|+|Y|\le a+b,
\]
we have
\[
    |X|\le a+b-|Y|.
\]
Hence
\[
\begin{aligned}
    H_X
    &\le |X|-x^*\\
    &<
    a+b-|Y|
    -
    \left(
        b-a+\left(1-\frac{2}{\alpha}\right)y^*
    \right)\\
    &=
    2a-|Y|-\left(1-\frac{2}{\alpha}\right)y^*\\
    &\le 2a.
\end{aligned}
\]
Thus
\[
    H_X\le 2|A|.
\]

We now bound the contribution of high arrivals into $X$. For any high
$v\in X$ and any $y\in Y_{<v}$, by triangle inequality, for every $a'\in A$,
\[
    d(y,a')\le d(y,v)+d(v,a').
\]
Rearranging gives
\[
    s(v,y)\le s(y,a')+d(v,a').
\]
Averaging over $a'\in A$ gives
\[
    s(v,y)
    \le
    \frac{s(y,A)}{|A|}
    +
    \frac{d(v,A)}{|A|}.
\]
Summing over $y\in Y_{<v}$,
\[
    s(v,Y_{<v})
    \le
    \frac{s(Y_{<v},A)}{|A|}
    +
    \frac{|Y_{<v}|}{|A|}d(v,A).
\]

For a high arrival into $X$, the previous argument also implies
\[
    |Y_{<v}|
    \le
    \frac{\alpha}{\alpha-1}|A|.
\]
Indeed, from
\[
    a+x
    >
    b+\left(1-\frac{2}{\alpha}\right)y
\]
and
\[
    x+y\le a+b,
\]
we get
\[
    \left(2-\frac{2}{\alpha}\right)y<2a,
\]
and therefore
\[
    y<\frac{\alpha}{\alpha-1}a.
\]
Thus, for every high arrival $v\in X$,
\[
    s(v,Y_{<v})
    \le
    \frac{s(Y,A)}{|A|}
    +
    \frac{\alpha}{\alpha-1}d(v,A).
\]
Summing over all high arrivals into $X$, and using $H_X\le 2|A|$, gives
\[
\begin{aligned}
    \sum_{\substack{v\in X\\ v\text{ high}}} s(v,Y_{<v})
    &\le
    \frac{H_X}{|A|}s(Y,A)
    +
    \frac{\alpha}{\alpha-1}d(X,A)\\
    &\le
    2s(Y,A)
    +
    \frac{\alpha}{\alpha-1}d(X,A).
\end{aligned}
\]

\paragraph{High case into $Y$.}
The argument is symmetric. A high arrival into $Y$ is a vertex $v\in Y$ such
that
\[
    s(v,X_{<v})
    >
    \alpha\,d(v,A\cup X_{<v}).
\]
By the same reasoning, the number of high arrivals into $Y$ is at most
$2|B|$, and
\[
    \sum_{\substack{v\in Y\\ v\text{ high}}} s(v,X_{<v})
    \le
    2s(X,B)
    +
    \frac{\alpha}{\alpha-1}d(Y,B).
\]

\paragraph{Putting the bounds together.}
Every pair in $X\times Y$ is counted exactly once, namely when the later of
the two endpoints arrives. Therefore,
\[
    s(X,Y)
\]
is the sum of the low-case and high-case contributions above. Combining the
three bounds, we get
\[
\begin{aligned}
    s(X,Y)
    &\le
    \alpha\bigl(d(X,B)+d(Y,A)+d(X,Y)\bigr)\\
    &\quad
    +\frac{\alpha}{\alpha-1}\bigl(d(X,A)+d(Y,B)\bigr)
    +2s(Y,A)+2s(X,B).
\end{aligned}
\]
The first two groups of terms are contained in $I^{\OPT}$, so
\[
    \alpha\bigl(d(X,B)+d(Y,A)+d(X,Y)\bigr)
    \le
    \alpha I^{\OPT},
\]
and
\[
    \frac{\alpha}{\alpha-1}\bigl(d(X,A)+d(Y,B)\bigr)
    \le
    \frac{\alpha}{\alpha-1}I^{\OPT}.
\]
\tienlong{The above two terms can be combined and bounded from above by $ \max(\alpha, \alpha/(\alpha-1)) I^{\OPT} $}
By assumption,
\[
    s(X,B)+d(X,A)+d(Y,B)+s(Y,A)
    \le
    cI^{\OPT}.
\]
In particular,
\[
    s(X,B)+s(Y,A)\le cI^{\OPT}.
\]
Therefore,
\[
    2s(X,B)+2s(Y,A)\le 2cI^{\OPT}.
\]
Combining these inequalities yields
\[
    s(X,Y)
    \le
    \left(
        \alpha+\frac{\alpha}{\alpha-1}+2c
    \right)I^{\OPT}.
\]
This completes the proof.
\end{proof}
